Un dels candidats a forat negre més pròxims a la Terra és A0620-00, que està situat a uns 3\,500 anys llum. Es calcula que la massa d'aquest forat negre és de $2,2\times 10^{31}\ \mathrm{kg}$. Encara que A0620-00 no és visible, s'ha detectat una estrella que descriu cercles amb un període orbital de 0,33 dies al voltant d'un lloc on no es detecta cap altre astre.

\figura[0.3]{fig1.png}

\begin{apartat}{1}
Deduïu la fórmula per a obtenir el radi d'una òrbita circular a partir de les magnituds proporcionades. Utilitzeu aquesta fórmula per a calcular el radi de l'òrbita de l'estrella que es mou al voltant d'A0620-00.
\end{apartat}

\begin{apartat}{1}
Calculeu la velocitat lineal i l'acceleració centrípeta de l'estrella i representeu els dos vectors $\boldsymbol{v}$ i $\boldsymbol{a}_\mathrm{c}$ sobre una figura similar a la d'aquest problema.
\end{apartat}

\blocetiquetat{Dada:}{$G = 6,67\times 10^{-11}\ \mathrm{N\,m^2\,kg^{-2}}$.}
